\documentclass[11pt,a4paper]{article}
\usepackage[T1]{fontenc}
\usepackage{XCharter}
\usepackage[margin=25mm]{geometry}
\usepackage{microtype,amsmath,amssymb,booktabs}
\usepackage{xurl}
\usepackage[colorlinks=true,urlcolor=blue,linkcolor=blue]{hyperref}
\setlength{\emergencystretch}{2em}
\title{When Is a Risk Certificate Worth Its Sampling Cost?\\
\large A synthetic deepening of the recursive CVaR branch}
\author{Pathfinder supervisor-led research follow-up}
\date{28 September 2026}
\begin{document}
\maketitle

\begin{abstract}
We instantiate the recursive paper's conditional-risk obedience certificate
in a two-signal routing model with a paid simulator, a reward budget and a
bounded fallback. This gives a conditional expected-cost comparison absent
from the parent paper's application-free statement. A fixed simulation grid
shows useful deployment on a separated instance, persistent abstention near
the feasibility boundary, and a cost penalty for oversampling. Knowing the
safe alternative exactly reduces both the uncertainty and the access bill.
This is a synthetic worked application, not field validation, a new QP
generation, or an independently reviewed manuscript.
\end{abstract}

\section{Research question and lineage}

The parent is \emph{Priced Certificates for Conditional-CVaR Obedience},
recursive Q1P2, clarified version 2 of 28 September 2026. It combines an
inherited deterministic reward graph with a simultaneous realised-data
confidence box. Its open questions include pricing counterfactual access and
specifying an abstention fallback. This follow-up supplies those ingredients
in a deliberately small example. It does not add evidence to the matched
repeat-versus-recurse comparison, whose allocations are already complete.

The protocol was written before executing the simulations. Code and output
are in \path{experiments/2026-09-28-risk-deepening/}. The result file binds
the parent, protocol and script by SHA-256. Historical papers are unchanged.
All claims below are either direct algebra in this declared model or an
application of the parent's stated confidence guarantee. No publication
priority or exhaustive external literature claim is made.

\section{A declared routing model}

Signals \(s=0,1\) are equally likely. The fixed policy prescribes action
\(A\) at signal 0 and \(B\) at signal 1. At signal \(s\), the prescribed
route has loss \(2X_s\), with \(X_s\sim\mathrm{Bernoulli}(p_s)\).
The alternative route has constant loss 1. Thus each action is prescribed
somewhere, as required by the parent. The agent evaluates conditional loss
using upper-tail CVaR with tail mass \(\beta=1/2\), and subtracts its
state-independent action reward. Recommendation-following breaks utility
ties. This is weak obedience, not unique or strict implementation.

Write \(c_s=\min\{4p_s,2\}\) for the prescribed route's true CVaR and
\(d=r(A)-r(B)\). Exact obedience is simply
\[
 c_0-1\leq d\leq 1-c_1.
\]
It is feasible exactly when \(c_0+c_1\leq2\). Among nonnegative rewards,
choose the point of this interval closest to zero. Its absolute value is
the oracle cap \(R^*\), and its expected payment is \(R^*/2\).

The planner pays operating cost \(g=0.2\) on deployment and expected
action reward \(|d|/2\). For evaluation we conservatively add a surcharge
\(J=2\) if \emph{any} true obedience constraint fails, rather than
modelling the deviating agent's route. Reward cap is \(R=0.6\).
Fallback costs \(F=1\), needs no agent obedience, and is always available.
Consequently downstream evaluation cost is at most
\(M=g+R/2+J=2.5\). These costs are stipulated in a common unit; they are
not derived from the agent's risk objective or measured operational losses.

Sampling uses a conditional simulator with price \(\kappa\) per loss draw.
This assumes access to deviation outcomes. On-policy logs alone do not
justify this simulator. Two access regimes are compared:
\begin{itemize}
\item Uniform: sample all four signal-action cells, charging \(4m\) draws.
  The parent's generic certificate treats every cell as uncertain, even
  though safe-cell observations in this example are constant.
\item Known safe: the alternative's constant loss is known from the model
  or an enforceable service contract. Sample only the two risky cells,
  charging \(2m\) draws. This assumption is stronger than observing a
  constant value in a finite batch.
\end{itemize}

\section{Certificate, fallback and an expected-cost bound}

For \(N=4\) or 2 respectively, the parent gives simultaneous radius
\[
 e=4\sqrt{\frac{\log(2N/\delta)}{2m}},\qquad \delta=0.05.
\]
The empirical risky CVaR is exactly \(\widehat c_s=\min(4K_s/m,2)\),
where \(K_s\) is a binomial count. Put \(t=2e\) in the uniform regime
and \(t=e\) in the known-safe regime. Since sampled safe losses equal 1,
the robust interval in this particular application is
\[
 \widehat c_0-1+t\leq d\leq 1-\widehat c_1-t.
\]
If this interval intersects \([-R,R]\), deploy using its point closest
to zero. If \(\widehat c_0+\widehat c_1-2t>2\), record a refutation
and use fallback. Otherwise use fallback, distinguishing an uncertain
budget from ordinary abstention. A failed robust budget check does not
certify failure of the true budget.

On simultaneous coverage, deployment is obedient and a refutation is sound.
This is not a guarantee on every batch. Suppose
\(\gamma=(2-c_0-c_1)/2>0\). Define
\[
 I=\begin{cases}4e&\text{uniform},\\2e&\text{known safe}.\end{cases}
\]
If \(I<\gamma\) and \(R^*+I\leq R\), every covered batch deploys
and its cap is at most \(R^*+I\). To see this, each robust difference
constraint exceeds its true counterpart by at most \(I\); their sum
is at most \(-2\gamma+2I<0\). The two-action path potential increases
by at most \(I\), and the budget condition permits that potential.
The uniform bound deliberately retains the parent's generic error-box
argument; the deterministic safe observations could also permit sharpening.

Let \(q=4m\) or \(2m\). Since downstream cost on coverage is at most
\(B=g+(R^*+I)/2\leq0.5<M\), splitting by the coverage event gives
\[
 \mathbb E[\text{total cost}]
 \leq \kappa q+(1-\delta)\left(g+\frac{R^*+I}{2}\right)+\delta M.
 \tag{1}
\]
This compares with the feasible prediction-free benchmark of always
falling back at cost 1. It is not a ratio against an unknown optimal system.
It also differs from the oracle-with-free-information benchmark, which
can deploy at \(g+R^*/2\) whenever feasible and within budget.

For \((p_0,p_1)=(0.3,0.1)\), we have \(R^*=0.2\) and
\(\gamma=0.2\). Known-safe access at \(m=4000\) and
\(\kappa=10^{-5}\) satisfies the premises and gives (1) approximately
\(0.579<1\). Uniform access at \(m=64000\) gives approximately
\(3.018\), not a useful guarantee at that price; at \(\kappa=10^{-6}\)
the same uniform bound is \(0.714<1\). These are sufficient bounds,
not optimal access-complexity results.

\section{Fixed-batch numerical experiment}

The protocol fixes five probability pairs, six batch sizes from 250 through
256000, two access regimes and two sampling prices. Each scenario-size
combination has 2000 independent batches, pairing regimes on the same risky
observations. There are 60000 batch pairs and 120000 regime decisions.
Binomial counts are sampled directly using Python 3.13.0 with seed 20260928.
No normal approximation is used. Guarantees apply per fixed trial, not
simultaneously across the grid; no adaptive stopping is claimed.

Table~\ref{tab:paid} shows the paid-reward instance at the higher sampling
price. Mean costs include fallback and the stipulated error surcharge.
The baseline without sampling costs exactly 1.

\begin{table}[ht]
\centering\small
\begin{tabular}{r l r r r}
\toprule
\(m\)&Access&Deployment (\%)&Mean cost&MC standard error\\
\midrule
1000&Known safe&63.2&0.629&0.00676\\
4000&Uniform&47.3&0.872&0.00681\\
4000&Known safe&100.0&0.428&0.00103\\
16000&Uniform&100.0&0.990&0.00016\\
16000&Known safe&100.0&0.644&0.00100\\
64000&Uniform&100.0&2.885&0.00008\\
64000&Known safe&100.0&1.593&0.00100\\
\bottomrule
\end{tabular}
\caption{Paid instance \((0.3,0.1)\), price \(10^{-5}\) per draw.
Deployment frequency is not correctness frequency.}\label{tab:paid}
\end{table}

Across the entire grid, six deployed rewards violate a true obedience
constraint. They all occur in the paid scenario with known-safe access:
one each at \(m=1000,4000,16000,64000\), and two at 256000.
All occur outside the recorded confidence box. There are 106 coverage
failures across the paired regime decisions, and no incorrect refutations.
The maximum observed false-decision frequency in a grid cell is 0.1\%;
neither that frequency nor zero errors in other cells proves coverage.
The surcharge is retained in every reported cost, not removed as an outlier.

The free-reward scenario \((0.1,0.1)\) readily deploys without payments.
The near-boundary scenario \((0.35,0.149)\) and exact boundary
\((0.35,0.15)\) abstain in every tested batch under both regimes.
For the infeasible scenario \((0.35,0.2)\), known-safe access refutes in
61.05\% of batches at \(m=4000\) and all batches at 16000; uniform
access refutes in 48\% at 16000 and all batches at 64000.
Refutation is informative but cannot repay sampling in this one-decision
model: fallback is still required, and could have been chosen for free.

\section{What this deepening adds}

First, specifying a fallback and a failure-path cap closes the parent's
missing expected-cost calculation, for this model only. The guarantee can
be economically useful: equation (1) is below fallback cost for a declared
instance, batch and sampling price.

Second, the access model materially changes the result. Known-safe access
both halves the charged draws and removes uncertainty from deviation loss.
These effects are combined here, not causally separated. It would be
misleading to describe the gain as a uniformly better algorithm on equal
information assumptions.

Third, more evidence is not always better for the planner. Once deployment
is already reliable, the sampling bill grows linearly while the reward
premium decreases only as the inverse square root of batch size. The grid
illustrates this tradeoff; selecting the best observed grid cell would be
post-selection, not a validated adaptive policy.

Finally, a small margin changes the practical question. At
\((0.35,0.149)\), \(\gamma=0.002\). The known-safe sufficient margin
condition alone requires more than 35 million samples per risky cell;
the uniform version requires more than 162 million per cell. At either
declared price those sufficient sample bills already exceed the fallback.
This does not prove that every valid certificate needs so many samples:
it diagnoses the conservatism and economic limitation of this particular
sufficient rule. A useful next investigation is a budget-aware sequential
certificate with valid time-uniform coverage, or amortising one certificate
over explicitly stationary repeated decisions. Neither is implemented here.

\section{Verification and limits}

Tests cover exact CVaR examples, interval signs, budget fallback,
confidence-box corners, access accounting and the sufficient deployment
premises. The simulation checks every definitive outcome against the true
array, including failures. These checks supplement the algebra above.
Results depend on a correct conditional simulator, stationary losses, known
safe loss in the stronger regime, and the chosen cost model. There is no
measured simulator cost, multi-type mechanism, report incentive analysis,
or field experiment. No external review was performed.

This is \emph{deepening}: pursuing an intermediate paper's application and
limitations. It is not another recombination and should not be counted as
a new success in the historical repeat-versus-recurse comparison.
\end{document}
